
\documentclass{article}
\usepackage{colortbl}
\usepackage{makecell}
\usepackage{multirow}
\usepackage{supertabular}

\begin{document}

\newcounter{utterance}

\centering \large Interaction Transcript for game `eqbench', experiment `eq\_bench', episode 15 with reflexion{-}llama{-}3.1{-}8b{-}instruct.
\vspace{24pt}

{ \footnotesize  \setcounter{utterance}{1}
\setlength{\tabcolsep}{0pt}
\begin{supertabular}{c@{$\;$}|p{.15\linewidth}@{}p{.15\linewidth}p{.15\linewidth}p{.15\linewidth}p{.15\linewidth}p{.15\linewidth}}
    \# & \multicolumn{2}{c}{Player} && \multicolumn{2}{c}{Game Master} \\
    \hline

    \theutterance \stepcounter{utterance}  
    & & & \multicolumn{4}{p{0.6\linewidth}}{
        \cellcolor[rgb]{0.9,0.9,0.9}{
            \makecell[{{p{\linewidth}}}]{
                \texttt{\tiny{[P1$\langle$GM]}}
                \texttt{You have attempted this same task before and did not succeed. These are the plans you wrote after those attempts.} \\
\texttt{Trial \#1: New plan: I was too focused on scoring high on the emotions that I thought were likely, but I didn't consider the context of the dialogue. Grayson's accusations are quite serious, and Dr. Sinclair's response suggests that he's not taking them lightly. I should have given more weight to the emotions that are more closely related to the situation. Next time, I will consider the tone and context of the dialogue more carefully and score the emotions based on how well they fit the situation. I will also make sure to give more thought to the intensity of the emotions, rather than just assigning a score based on how likely they seem.} \\
\\ 
\texttt{Trial \#2: New plan: In my previous attempt, I focused too much on the emotions that I thought were likely, but I didn't consider the context of the dialogue. Grayson's accusations are quite serious, and Dr. Sinclair's response suggests that he's not taking them lightly. I should have given more weight to the emotions that are more closely related to the situation. Next time, I will consider the tone and context of the dialogue more carefully and score the emotions based on how well they fit the situation. I will also make sure to give more thought to the intensity of the emotions, rather than just assigning a score based on how likely they seem.} \\
\\ 
\texttt{To achieve this, I will analyze the dialogue more closely and identify the key elements that contribute to the emotional tone. I will then use this analysis to assign scores to each emotion, taking into account the context and tone of the dialogue. For example, I will consider how Dr. Sinclair's response suggests that he's not taking Grayson's accusations lightly, which would suggest a higher score for emotions like Offended or Skeptical. I will also consider how Grayson's accusations are quite serious, which would suggest a higher score for emotions like Skeptical or Offended.} \\
\\ 
\texttt{Here is my revised plan:} \\
\\ 
\texttt{Surprised: 2} \\
\texttt{Amused: 0} \\
\texttt{Offended: 9} \\
\texttt{Skeptical: 7} \\
\\ 
\texttt{Trial \#3: New plan: In my previous attempts, I focused too much on the emotions that I thought were likely, but I didn't consider the context of the dialogue carefully enough. I should have given more weight to the emotions that are more closely related to the situation. Next time, I will analyze the dialogue more closely and identify the key elements that contribute to the emotional tone. I will then use this analysis to assign scores to each emotion, taking into account the context and tone of the dialogue.} \\
\\ 
\texttt{To achieve this, I will pay closer attention to the language and tone used by the characters, as well as the situation and context of the dialogue. I will also consider how the characters' responses and actions contribute to the emotional tone, and use this information to assign scores to each emotion.} \\
\\ 
\texttt{Specifically, I will look for clues such as:} \\
\\ 
\texttt{* The tone of the characters' language (e.g. sarcastic, serious, etc.)} \\
\texttt{* The level of aggression or defensiveness in the characters' responses} \\
\texttt{* The context of the situation (e.g. a serious accusation, a personal attack, etc.)} \\
\texttt{* The characters' body language and nonverbal cues (e.g. facial expressions, tone of voice, etc.)} \\
\\ 
\texttt{By considering these factors, I will be able to assign scores to each emotion that are more accurate and reflective of the emotional tone of the dialogue.} \\
\\ 
\texttt{Surprised: 0} \\
\texttt{Amused: 0} \\
\texttt{Offended: 10} \\
\texttt{Skeptical: 10} \\
\\ 
\\ 
\\ 
\texttt{Your task is to predict the likely emotional responses of a character in this dialogue:} \\
\\ 
\texttt{Grayson: You know, Doc, I've been watching your work for years. All that time and money spent on Mars, and what have you got to show for it? } \\
\texttt{Dr. Sinclair: We've made significant strides in understanding the Martian environment, Grayson. It's not about instant results. } \\
\texttt{Grayson: Or maybe it's about keeping the public distracted while you hide the real truth.} \\
\texttt{Dr. Sinclair: And what truth would that be, Mr. Hawk?} \\
\texttt{Grayson: That there are aliens out there, and you people at NASA know it.} \\
\texttt{{[}End dialogue{]}} \\
\\ 
\texttt{At the end of this dialogue, Dr. Sinclair would feel...} \\
\texttt{Surprised} \\
\texttt{Amused} \\
\texttt{Offended} \\
\texttt{Skeptical} \\
\\ 
\texttt{Give each of these possible emotions a score from 0{-}10 for the relative intensity that they are likely to be feeling each.} \\
\\ 
\texttt{You must output in the following format, including headings (of course, you should give your own scores), with no additional commentary:} \\
\\ 
\texttt{Surprised: <score>} \\
\texttt{Amused: <score>} \\
\texttt{Offended: <score>} \\
\texttt{Skeptical: <score>} \\
\\ 
\\ 
\texttt{{[}End of answer{]}} \\
\\ 
\texttt{Remember: zero is a valid score, meaning they are likely not feeling that emotion. You must score at least one emotion > 0.} \\
\\ 
\texttt{Your answer:} \\
            }
        }
    }
     \\ \\

    \theutterance \stepcounter{utterance}  
    & \multicolumn{4}{p{0.6\linewidth}}{
        \cellcolor[rgb]{0.561,0.737,0.561}{
            \makecell[{{p{\linewidth}}}]{
                \texttt{\tiny{[P1$\rangle$GM]}}
                \texttt{Surprised: 0} \\
\texttt{Amused: 0} \\
\texttt{Offended: 8} \\
\texttt{Skeptical: 9} \\
            }
        }
    }
    & & \\ \\

    \theutterance \stepcounter{utterance}  
    & & & \multicolumn{2}{p{0.3\linewidth}}{
        \cellcolor[rgb]{0.9,0.9,0.9}{
            \makecell[{{p{\linewidth}}}]{
                \texttt{\tiny{[GM$|$GM]}}
                \texttt{Surprised: 0} \\
\texttt{Amused: 0} \\
\texttt{Offended: 8} \\
\texttt{Skeptical: 9} \\
            }
        }
    }
    & & \\ \\

    \theutterance \stepcounter{utterance}  
    & & & \multicolumn{2}{p{0.3\linewidth}}{
        \cellcolor[rgb]{0.9,0.9,0.9}{
            \makecell[{{p{\linewidth}}}]{
                \texttt{\tiny{[GM$|$GM]}}
                \texttt{Surprised: 4} \\
\texttt{Amused: 3} \\
\texttt{Offended: 0} \\
\texttt{Skeptical: 8} \\
            }
        }
    }
    & & \\ \\

    \theutterance \stepcounter{utterance}  
    & & & \multicolumn{2}{p{0.3\linewidth}}{
        \cellcolor[rgb]{0.9,0.9,0.9}{
            \makecell[{{p{\linewidth}}}]{
                \texttt{\tiny{[GM$|$GM]}}
                \texttt{game\_result = LOSE} \\
            }
        }
    }
    & & \\ \\

    \theutterance \stepcounter{utterance}  
    & & \multicolumn{4}{p{0.6\linewidth}}{
        \cellcolor[rgb]{0.95,0.95,0.95}{
            \makecell[{{p{\linewidth}}}]{
                \texttt{\tiny{[P1$\leftrightarrow$P]}}
                \texttt{Post{-}episode Reflexion (not sent to GM):} \\
\texttt{New plan: In my previous attempts, I focused too much on the emotions that I thought were likely, but I didn't consider the context of the dialogue carefully enough. I should have given more weight to the emotions that are more closely related to the situation. Next time, I will analyze the dialogue more closely and identify the key elements that contribute to the emotional tone. I will then use this analysis to assign scores to each emotion, taking into account the context and tone of the dialogue.} \\
\\ 
\texttt{To achieve this, I will pay closer attention to the language and tone used by the characters, as well as the situation and context of the dialogue. I will also consider how the characters' responses and actions contribute to the emotional tone, and use this information to assign scores to each emotion.} \\
\\ 
\texttt{Specifically, I will look for clues such as:} \\
\\ 
\texttt{* The tone of the characters' language (e.g. sarcastic, serious, etc.)} \\
\texttt{* The level of aggression or defensiveness in the characters' responses} \\
\texttt{* The context of the situation (e.g. a serious accusation, a personal attack, etc.)} \\
\texttt{* The characters' body language and nonverbal cues (e.g. facial expressions, tone of voice, etc.)} \\
\\ 
\texttt{By considering these factors, I will be able to assign scores to each emotion that are more accurate and reflective of the emotional tone of the dialogue.} \\
\\ 
\texttt{I will also make sure to consider the intensity of the emotions, rather than just assigning a score based on how likely they seem. For example, Grayson's accusations are quite serious, which would suggest a higher score for emotions like Skeptical or Offended.} \\
\\ 
\texttt{Surprised: 0} \\
\texttt{Amused: 0} \\
\texttt{Offended: 10} \\
\texttt{Skeptical: 10} \\
            }
        }
    }
    & \\ \\

\end{supertabular}
}

\end{document}
